\section{Preliminaries: Basics of matrix analysis} \label{subsec:Matrix-analysis}


We begin this chapter by gathering a few elementary materials in matrix analysis that prove useful for our theoretical development. The readers familiar with matrix analysis can proceed directly to Section~\ref{sec:preliminary-distance-angles}. 
 

 
\paragraph{Unitarily invariant norms.}



Among all matrix norms, the family of unitarily invariant norms defined below is of central interest, which subsumes  as special cases the spectral norm $\|\cdot\|$ and the Frobenius norm $\|\cdot\|_{\mathrm{F}}$.
%
\begin{definition}
A matrix norm $\vertiii{\cdot}$ on $\mathbb{R}^{m\times n}$ is said to be unitarily invariant if
%
\[
	\vertiii{\bm{A}}=\vertiiibig{\bm{U}^{\top}\bm{A}\bm{V}}
\]
%
holds for any matrix $\bm{A}\in\mathbb{R}^{m\times n}$ and any two square orthonormal
matrices $\bm{U}\in\mathcal{O}^{m\times m}$ and $\bm{V}\in\mathcal{O}^{n\times n}$.
\end{definition}
%
This class of matrix norms enjoys several useful properties, as summarized in the following lemma. The proof can be found in \citet[Theorem 3.9]{stewart1990matrix}.
%
\begin{lemma}
\label{prop:unitary_norm_relation}
For any unitarily invariant norm $\vertiii{\cdot}$, one has 
%
\begin{align*}
\vertiii{\bm{A}\bm{B}} & \leq\vertiii{\bm{A}} \cdot  \left\Vert \bm{B}\right\Vert, &  & \vertiii{\bm{A}\bm{B}}\leq\vertiii{\bm{B}} \cdot \left\Vert \bm{A}\right\Vert,\\
\vertiii{\bm{A}\bm{B}} & \geq\vertiii{\bm{A}} \, \sigma_{\min}\left(\bm{B}\right), &  & \vertiii{\bm{A}\bm{B}}\geq\vertiii{\bm{B}} \, \sigma_{\min}\left(\bm{A}\right).
\end{align*}
%
%where $\sigma_{\min}(\bm{A})$ denotes the smallest singular value of a matrix $\bm{A}$. 
\end{lemma}



	
\paragraph{Perturbation bounds for eigenvalues and singular values.}

Next, we review classical perturbation
bounds for eigenvalues of symmetric matrices and for singular values of general matrices. 

\begin{lemma}[Weyl's inequality for eigenvalues]
\label{lemma:weyl}
Let $\bm{A},\bm{E}\in\mathbb{R}^{n\times n}$ be two real symmetric matrices.
	For every $1\leq i\leq n$, the $i$-th largest eigenvalues of $\bm{A}$ and $\bm{A}+\bm{E}$ obey
%
\begin{equation}
	\left|\lambda_{i}\left(\bm{A}\right)-\lambda_{i}\left(\bm{A} +\bm{E}\right)\right|\leq\left\Vert \bm{E}\right\Vert .
	\label{eq:weyl}
\end{equation}
%
\end{lemma}
\begin{proof}See Equation (1.63) in \citet{Tao2012RMT}. \end{proof}



\begin{lemma}[Weyl's inequality for singular values]
\label{lemma:weyls-singular-value}
Let $\bm{A},\bm{E}\in\mathbb{R}^{m\times n}$ be two general matrices.
Then for every $1\leq i\leq \min\{m,n\}$, the $i$-th largest singular values of $\bm{A}$ and $\bm{A}+\bm{E}$ obey 
%
\[
	\left|\sigma_{i}\left(\bm{A}+\bm{E}\right)-\sigma_{i}\left(\bm{A}\right)\right|\leq\left\Vert \bm{E}\right\Vert .
\]
%
\end{lemma}
\begin{proof}See Exercise 1.3.22 in \citet{Tao2012RMT}. \end{proof}



An immediate implication of Lemma~\ref{lemma:weyl} (resp.~Lemma~\ref{lemma:weyls-singular-value}) is that the eigenvalues of a real symmetric matrix (resp.~the singular values of a general matrix) 
 are stable vis-\`a-vis small perturbations. 




